\documentclass[11pt,a4paper]{article}
\usepackage[utf8]{inputenc}
\usepackage[T1]{fontenc}
\usepackage[spanish,es-nodecimaldot]{babel}
\usepackage[margin=2.2cm]{geometry}
\usepackage{amsmath,amssymb,booktabs,longtable,tabularx,array}
\usepackage{enumitem,xcolor,hyperref,xurl,tikz}
\usetikzlibrary{arrows.meta,positioning}
\hypersetup{colorlinks=true,linkcolor=blue!55!black,urlcolor=blue!55!black}
\setlist{leftmargin=*,itemsep=3pt}
\setlength{\parskip}{5pt}
\setlength{\emergencystretch}{1.5em}
\newcommand{\src}[1]{\nolinkurl{#1}}
\newcommand{\dd}{\mathrm d}
\newcommand{\unit}[1]{\mathrm{#1}}
\newcommand{\warn}[1]{\par\noindent\colorbox{blue!8}{\parbox{.965\linewidth}{\textbf{Lectura crítica.} #1}}\par}
\title{Simulación de inundación DG2 por región de interés\\\large Ecuaciones, discretización, entradas y ciclo de vida verificable}
\author{TerraPulse AI · Documento técnico independiente}
\date{Versión documental de la implementación observada en la rama simulacion-inundancion}
\begin{document}
\maketitle
\begin{abstract}
Se describe la ruta implementada que, a petición autenticada, transforma una región de interés (ROI) y un dominio declarados en un escenario de lluvia \emph{sintética}, prepara terreno NASADEM, ejecuta LISFLOOD-FP 8.0 DG2 CPU y ofrece mapas numéricos de profundidad. Se distinguen tres niveles: el modelo físico de aguas someras, su aproximación numérica DG2 y el adaptador de archivos, procesos y API. Cada ecuación identifica magnitudes y unidades; las afirmaciones de ejecución se vinculan al código. Un resultado integrable y descargable no es una predicción validada de inundación: todavía faltan una elección científicamente defendible del dominio, caudales, rugosidad, infiltración, balance comprobado y contraste independiente. No se atribuye a NISAR ninguna profundidad, lluvia ni causa de cambio.
\end{abstract}
\tableofcontents\newpage

\section{Guía de lectura y alcance}
\label{sec:alcance}
\begin{center}
\begin{tikzpicture}[node distance=3mm,box/.style={draw,rounded corners,fill=blue!6,align=center,text width=3.15cm,minimum height=1.3cm},arr/.style={-{Stealth},thick}]
\node[box] (a) {Modelo continuo\\SWE 2D, gravedad, lecho, Manning};
\node[box,right=of a] (b) {Método numérico\\coeficientes DG2, HLL, RK2, CFL};
\node[box,right=of b] (c) {Adaptador operativo\\GeoTIFF / NetCDF, proceso CPU, API/COG};
\draw[arr] (a)--(b); \draw[arr] (b)--(c);
\end{tikzpicture}
\end{center}
\noindent\textbf{Pregunta y respuesta.} ¿Qué calcula la cadena disponible? La profundidad $h$ de una lámina superficial, en metros \emph{sobre el lecho local}, para un dominio métrico rectangular, bajo lluvia hipotética espacialmente uniforme en las celdas seleccionadas por la ROI, estado inicial seco, Manning constante y borde idealizado elegido. Los mapas muestran profundidades de cuadros y máximo; no son mapas medidos por satélite. El cálculo no proporciona al cliente velocidad, tiempo de llegada ni extensión observada. El usuario elige entradas y supuestos; el servicio no determina si son físicamente correctos.

El escenario se limita a terreno declarado no urbano. \texttt{domain\_nonurban:true} expresa una declaración del solicitante, no una clasificación automática de cobertura del suelo. Una ROI de consulta no delimita necesariamente la cuenca que alimenta el flujo; el dominio debe contemplar aportes y salidas. Aquí no se representan calles, alcantarillas, edificios, batimetría detallada ni flujo subterráneo. Si esas omisiones controlan el evento, no interpretar el resultado como una simulación regional válida. La ruta ROI actual es \src{pipeline/hydrology/roi_runner.py}; los casos registrados y \texttt{prepare/} son contratos diferentes y no se mezclan.

\warn{La verificación de formato y de ejecución no equivale a validación hidráulica. Los valores numéricos de una solicitud de ejemplo más adelante son únicamente ilustrativos; no son recomendaciones ni resultados observados.}

\section{Magnitudes, referencias y supuestos físicos}
\label{sec:simbolos}
Se suponen coordenadas horizontales proyectadas $x,y$ en metros (no grados), tiempo $t$ en segundos desde el comienzo \emph{simulado} y lecho estacionario $z_b(x,y)$ en metros de elevación respecto de un datum vertical declarado. $h(x,y,t)\geq0$ es altura de agua \emph{sobre ese lecho}, no altitud absoluta. La superficie libre $\eta=z_b+h$ es elevación respecto del \emph{mismo} datum; restar cotas de datum diferente no produce una profundidad. NASADEM\_HGT aporta elevación referida al geoide EGM96; la reproyección horizontal no efectúa una conversión vertical. $u,v$ son componentes de velocidad horizontal \emph{promediadas en la profundidad}, en m/s; $g$ es aceleración gravitatoria en m/s$^2$.

Definimos $q_x=hu$ y $q_y=hv$ (m$^2$/s): cada uno es volumen que cruza una sección vertical \emph{por segundo y por metro de ancho perpendicular}. En cambio $Q=\int_B\boldsymbol q\cdot\boldsymbol n_{\!f}\,\dd\ell$ tiene m$^3$/s a través del ancho $B$ de una frontera/sección; $\boldsymbol n_{\!f}$ es la normal horizontal unitaria, no la rugosidad de Manning $n_M$. Para flujo uniforme sobre ancho $B$, $Q=q_{\perp}B$, pero no para una sección heterogénea. $Q/A_c$ (m/s) divide por área \emph{transversal mojada} $A_c$ (m$^2$), velocidad media; $Q/A_p$ (m/s) divide por área \emph{horizontal en planta} $A_p$ (m$^2$), tasa areal equivalente únicamente si se distribuye allí el volumen. El símbolo $q_s$ designará exclusivamente otra \emph{fuente areal} de profundidad (m/s), nunca $q_x$ o $Q$. No hay entrada $q_s$ general en esta API.

\begin{longtable}{p{.18\linewidth}p{.17\linewidth}p{.57\linewidth}}
\toprule Símbolo & Unidad SI & Referencia, denominador o interpretación\\\midrule\endhead
$h,z_b,\eta$ & m & Profundidad desde lecho; cota de lecho y cota de superficie libre desde datum común. $\eta-z_b=h$.\\
$x,y,\Delta x,\Delta y$ & m & Posición horizontal y longitud de lado de celda proyectada; las diferencias no se miden en grados de latitud.\\
$t,\Delta t$ & s & Tiempo desde origen simulado y duración de un paso; sin fecha UTC inferida.\\
$u,v,c=\sqrt{gh}$ & m/s & Velocidad media de columna y celeridad gravitatoria; $gh$ tiene m$^2$/s$^2$.\\
$q_x,q_y$ & m$^2$/s & Caudal \emph{por unidad de ancho} de sección vertical, no m$^3$/s ni profundidad por tiempo.\\
$Q$ & m$^3$/s & Volumen por segundo a través de sección completa.\\
$r,i,q_s$ & m/s & Lámina ganada por lluvia, perdida por infiltración y fuente neta por \emph{área horizontal y tiempo}: m$^3$/(m$^2$ s). $i,q_s$ son conceptuales, no aplicados por la ruta ROI.\\
$n_M$ & s/m$^{1/3}$ & Coeficiente Manning en la convención SI del solver, no una normal ni un conteo; valor requerido explícito en solicitud.\\
$S_{f,x},S_{f,y}$ & 1 & Pérdida de altura de energía por longitud horizontal (m/m), pendiente de fricción con signo de velocidad.\\
$\mathcal F,\mathcal G$ & (m$^2$/s, m$^3$/s$^2$, m$^3$/s$^2$) & Primera componente flujo de masa/profundidad por ancho; otras dos, flujos de momento integrado en la columna por longitud de borde y tiempo.\\
$\Delta A=\Delta x\Delta y$ & m$^2$ & Área horizontal de una celda, usada para convertir lámina a volumen.\\\bottomrule
\end{longtable}

El modelo de aguas someras integra verticalmente, supone presión aproximadamente hidrostática y representa escala horizontal mucho mayor que la profundidad característica. No resuelve velocidad vertical ni turbulencia local 3D. El lecho puede ser especialmente incierto bajo agua, vegetación, puentes y cauces estrechos: refinar una malla remuestreada no crea topografía ausente.

\section{Ecuaciones continuas y conservación}
\label{sec:ecuaciones}
Para $\boldsymbol U=(h,q_x,q_y)^{\mathrm T}$, la forma conservativa bidimensional del modelo base es
\begin{align}
\partial_t\boldsymbol U+\partial_x\mathcal F(\boldsymbol U)+\partial_y\mathcal G(\boldsymbol U)&=\boldsymbol S_b+\boldsymbol S_f+\boldsymbol S_r,\label{eq:conserv}\\
\mathcal F&=\begin{pmatrix}q_x\\ q_x^2/h+gh^2/2\\q_xq_y/h\end{pmatrix},\qquad
\mathcal G=\begin{pmatrix}q_y\\q_xq_y/h\\q_y^2/h+gh^2/2\end{pmatrix},\label{eq:flux}\\
\boldsymbol S_b&=\begin{pmatrix}0\\-gh\partial_xz_b\\-gh\partial_yz_b\end{pmatrix},\quad
\boldsymbol S_r=\begin{pmatrix}r\\0\\0\end{pmatrix},\quad
\boldsymbol S_f=\begin{pmatrix}0\\-ghS_{f,x}\\-ghS_{f,y}\end{pmatrix}.\label{eq:sources}
\end{align}
En $h=0$, los cocientes en \eqref{eq:flux} no se evalúan directamente: el código aplica umbrales y reconstrucción de caras. La escritura continua representa la física \emph{idealizada}; no reproduce por sí misma la secuencia discreta de operadores. La lluvia se añade a profundidad y no aporta momento explícito en la ruta DG2 consultada. Generalmente podría escribirse $\partial_t h+\nabla\cdot\boldsymbol q=r-i+q_s$, con términos de momento asociados a intercambios de masa si corresponde; \textbf{aquí} no se invocan infiltración ni fuente areal adicional mediante la API, por lo que en el interior el aporte externo implementado es $r$. No atribuir soporte físico a opciones que el parser nativo pudiera reconocer sin que el bucle DG2 las aplique.

\subsection{Comprobación dimensional término por término}
En continuidad, $[\partial_t h]=\unit{m/s}$ y $[\partial_xq_x]=[\partial_yq_y]=(\unit{m^2/s})/\unit m=\unit{m/s}$. La precipitación $r$ se define por \emph{área de planta}, de modo que $r\Delta A\Delta t$ es m$^3$. Ninguna intensidad en mm/h puede sumarse directamente a $h$ en m sin convertirla y multiplicarla por segundos.

En cada componente de momento, $[\partial_t q_x]=\unit{m^2/s^2}$. $[q_x^2/h]=(\unit{m^2/s})^2/\unit m=\unit{m^3/s^2}$ y $[gh^2]=\unit{m^3/s^2}$; las derivadas $\partial_x,\partial_y$ dividen por distancia horizontal, dando m$^2$/s$^2$. El gradiente $\partial_xz_b$ compara metros verticales entre metros horizontales (adimensional); por tanto $[gh\partial_xz_b]=\unit{m^2/s^2}$. También $[gh S_{f,x}]=\unit{m^2/s^2}$ porque $S_{f,x}$ es pendiente adimensional. Un flujo de momento que cruza una cara, multiplicado por longitud de esa cara, tiene m$^4$/s$^2$; al dividir por el área de celda m$^2$ vuelve a m$^2$/s$^2$. Las identidades dimensionales no prueban ni exactitud de DEM ni conservación del programa.

\subsection{Manning, pendiente y operador efectivamente aplicado}
Con $w=\sqrt{u^2+v^2}$ (m/s), la convención continua de pendiente de fricción es
\begin{equation}
 S_{f,x}=\frac{n_M^2uw}{h^{4/3}},\quad S_{f,y}=\frac{n_M^2vw}{h^{4/3}},\qquad
 \boldsymbol S_{f,\mathrm{mom}}=-\frac{g n_M^2}{h^{1/3}}(uw,vw)^{\mathrm T}.\label{eq:mann}
\end{equation}
\[
[n_M^2uw/h^{4/3}]=(\unit{s^2/m^{2/3}})(\unit{m^2/s^2})/\unit{m^{4/3}}=1.
\] El factor $C_f=gn_M^2/h^{1/3}$ es adimensional: \[
(\unit{m/s^2})(\unit{s^2/m^{2/3}})/\unit{m^{1/3}}=1.
\] $C_fuw$ tiene m$^2$/s$^2$, \emph{no} es la pendiente adimensional $S_{f,x}$. En el código \src{swe/dg2/friction.cpp} se llama \texttt{Sfx} al término de momento negativo, no a la pendiente de energía definida en \eqref{eq:mann}. Usar una sola nomenclatura sin esta distinción daría ecuaciones incompatibles.

Para las medias y, cuando aplica, valores en puntos de Gauss, el método nativo calcula
\begin{align}
 C_f&=gn_M^2/h^{1/3}, & D_x&=1+\Delta t\frac{C_f}{h}\frac{2u^2+v^2}{w},
 & D_y&=1+\Delta t\frac{C_f}{h}\frac{u^2+2v^2}{w},\label{eq:fricden}\\
 q_x^{\mathrm{nuevo}}&=q_x+\Delta t\frac{-C_fuw}{D_x}, &
 q_y^{\mathrm{nuevo}}&=q_y+\Delta t\frac{-C_fvw}{D_y}.\label{eq:fricdisc}
\end{align}
$[\Delta t(C_f/h)(u^2/w)]=\unit s(1/\unit m)(\unit{m/s})=1$, así que $D_x,D_y$ son adimensionales. Los términos añadidos a $q_x,q_y$ tienen $(\unit s)(\unit{m^2/s^2})=\unit{m^2/s}$. El código evita dividir cuando $h\leq$ \texttt{DepthThresh} o ambas componentes de velocidad son menores que \texttt{SpeedThresh}; actualiza también coeficientes de pendiente desde evaluaciones en Gauss. \textbf{Este operador se aplica fuera de las dos etapas RK}; la presencia de RK2 no demuestra segundo orden temporal global con lluvia y fricción partidas. Referencia: \src{runtime/source/swe/dg2/friction.cpp}, funciones \texttt{apply\_friction} y \texttt{friction}; en adelante se abrevia \texttt{runtime/source} como \texttt{fuente nativa}.

\section{De la ecuación a DG2}
\label{sec:dg2}
\subsection{Representación modal y terreno}
Sea $K_{ij}$ una celda de lado $\Delta x,\Delta y$, centro $(x_i,y_j)$ y coordenadas locales $\xi=2(x-x_i)/\Delta x$, $\zeta=2(y-y_j)/\Delta y$, ambas entre $-1$ y $1$ y \emph{adimensionales}. No se usa $\eta$ para referencia espacial: $\eta$ ya es la superficie libre. La aproximación planar para cada campo $W\in\{h,q_x,q_y,z_b\}$ es
\begin{equation}
 W_h(x,y)=W_0+\sqrt3 W_{1x}\xi+\sqrt3 W_{1y}\zeta.\label{eq:basis}
\end{equation}
$W_0$ representa la media de celda y $W_{1x},W_{1y}$ amplitudes modales con \emph{las mismas unidades que $W$}; no son $\partial_xW,\partial_yW$ en unidades por metro. En efecto $\partial_xW_h=2\sqrt3W_{1x}/\Delta x$. Esta base da valores en las caras $W_0\pm\sqrt3W_{1x}$ y en las posiciones de Gauss de un eje $W_0\pm W_{1x}$. Los índices $1x,1y$ se refieren a los grados de libertad, no a bandas satelitales. El código \src{swe/dg2/modifiedvars.cpp} define \texttt{limit\_neg/pos} y \texttt{gauss\_lower/upper} con esas combinaciones; \src{swe/dg2.cpp} actualiza los nueve coeficientes de $h,q_x,q_y$ en cada paso.

El preprocesador oficial \texttt{generateDG2DEM} toma el raster \texttt{terrain.dem.raw} como muestras de vértice (con índices en el borde acotados a la última celda). Si NW, NE, SW y SE son las cuatro muestras de un elemento, su proyección implementada es
\begin{equation}
 z_0=\frac{NW+NE+SW+SE}{4},\quad
 z_{1x}=\frac{NE+SE-NW-SW}{4\sqrt3},\quad
 z_{1y}=\frac{NE+NW-SE-SW}{4\sqrt3}.\label{eq:vertices}
\end{equation}
Si NW y NE son iguales y SW y SE menores, $z_{1y}>0$ al tomar $y$ creciente hacia el norte. Todas estas magnitudes tienen metros EGM96: dividir por $\Delta x$ o $\Delta y$ para formar un gradiente cambia la unidad. La media DG2 puede diferir del valor original de una muestra DEM; no declarar que el asistente conserva cada pixel sin cambios. El adaptador rechaza huecos y los centinelas $-9999$ y $10^{10}$; el parámetro \texttt{nodata\_elevation 10000000000} evita confundir elevación cero con nodata en la utilidad. Véanse \src{preprocess/generateDG2DEM.cpp}, \src{swe/dg2/fields.cpp} (\texttt{initialise\_field}) y \src{pipeline/hydrology/dem_dg2.py}.

\subsection{Residuos, caras y lecho}
La formulación DG proyecta el balance sobre $1,\sqrt3\xi,\sqrt3\zeta$. Esquemáticamente, para cada función base $\phi_a$ y componente conservada, la forma débil por celda es
\begin{equation}
 \int_K \phi_a\,\partial_t U_h\,\dd A =
 \int_K\!\bigl(\mathcal F(U_h)\partial_x\phi_a+\mathcal G(U_h)\partial_y\phi_a+\phi_a S_b\bigr)\dd A
 -\int_{\partial K}\!\phi_a\widehat{\mathcal H}(U_h^-,U_h^+;\boldsymbol n_{\!f})\dd\ell.
 \label{eq:weak}
\end{equation}
Esta es una \emph{explicación matemática} de proyección, no una transcripción literal de cada bucle; el código evalúa flujos en caras y Gauss, reconstruye el lecho y calcula \texttt{L0}, \texttt{L1x}, \texttt{L1y}. Para el modo medio de profundidad, la ecuación de actualización sin lluvia dentro de $L$ es
\begin{equation}
 L_0(h)=-\frac{\widehat F^h_E-\widehat F^h_W}{\Delta x}
          -\frac{\widehat G^h_N-\widehat G^h_S}{\Delta y}.\label{eq:mean}
\end{equation}
En $x$ o $y$, $[\widehat F^h]=\unit{m^2/s}$; al dividir por m el residuo tiene m/s. Los modos lineales no se obtienen copiando \eqref{eq:mean}: \texttt{L1x/L1y} incluyen flujos físicos en los dos puntos de Gauss, contribuciones de caras y el término de lecho correspondiente al momento. Cuando \texttt{thin\_depth} se cumple, los incrementos lineales se anulan en \texttt{L}. Por eso describir el ejecutable como un volumen finito de solo media sería inexacto. El lecho entra como $-gh\nabla z_b$ con reconstrucción propia, no mediante una resta de alturas en un PNG.

En una cara entre reconstrucciones izquierda y derecha, \src{swe/dg2/modifiedvars.cpp} calcula $z^*=\max(z_b^-,z_b^+)$ y, para cada lado, $h^*=\max(0,\eta^{\pm}-z^*)$; si la profundidad original es suficientemente húmeda, $q^*=h^*(q/h)$, en caso contrario $q^*=0$. Esta reconstrucción hidroestática da sentido a estados de cara secos y busca coherencia con el equilibrio de lago en reposo, pero \emph{no demuestra} conservación exacta, positividad global ni balance bien equilibrado en este caso sin ensayos. Los términos \texttt{bed\_source\_0x/0y/1x/1y} de \src{swe/dg2.cpp} acompañan los flujos en la actualización modal.

\subsection{HLL y componente transversal}
Para cada cara, se resuelve aproximadamente un problema normal 1D con los estados reconstruidos y las componentes de \eqref{eq:flux}. Si $S^-<0<S^+$, las dos primeras componentes del flujo HLL corresponden a
\begin{equation}
 \widehat F_{\mathrm{HLL}}^{(h,q_\perp)}=
 \frac{S^+F_L^{(h,q_\perp)}-S^-F_R^{(h,q_\perp)}
       +S^-S^+(U_R-U_L)^{(h,q_\perp)}}{S^+-S^-}.\label{eq:hll}
\end{equation}
Si $S^-\geq0$ se toma $F_L$; si $S^+<0$, $F_R$. $S^{\pm}$ son velocidades de ondas en m/s, acotadas en la fuente con $u\pm\sqrt{gh}$ y estimadores intermedios, con ramas especiales para estados secos. En la rama intermedia, \emph{la componente transversal no se calcula simplemente aplicando \eqref{eq:hll} a las tres componentes}: el código determina una onda de contacto intermedia y usa el flujo de profundidad central multiplicado por la velocidad transversal del lado elegido. La rutina en $y$ intercambia las componentes normales y transversales. Esto importa al comparar una ecuación HLL de libro con \src{swe/hll_include.h} y \src{swe/flux.cpp}. Chequeo: $[S^-S^+(h_R-h_L)]=\unit{(m/s)^2m}=\unit{m^3/s^2}$ y al dividir por $S^+-S^-$ (m/s) da m$^2$/s, como el flujo de profundidad; para momento, $(\unit{m/s})^2(\unit{m^2/s})/(\unit{m/s})=\unit{m^3/s^2}$.

\subsection{Dos etapas temporales y restricciones de estabilidad}
Al llamar $L$ al vector de residuos de los nueve coeficientes (con flujo y lecho), el núcleo ejecuta
\begin{equation}
 U^{(1)}=U^n+\Delta t L(U^n),\qquad
 U^{n+1}=\tfrac12\bigl(U^n+U^{(1)}+\Delta t L(U^{(1)})\bigr).\label{eq:rk}
\end{equation}
Es RK de dos etapas para \emph{ese} residuo: no se debe atribuir automáticamente segundo orden temporal al acoplamiento completo. \src{swe/dg2.cpp} aplica lluvia, actualiza condiciones, opcionalmente recalcula $\Delta t$, aplica fricción y luego \texttt{rk\_stage1/rk\_stage2}; a continuación actualiza máximos, acumula estadísticas de flujo, incrementa $t$ y escribe cuando se alcanzan los umbrales de guardado. No hay llamada a infiltración en \texttt{dg2::solve}. La aplicación de lluvia precede incluso al recálculo CFL: al iniciar un ciclo usa el \texttt{Tstep} vigente para $h\mathrel{+}=r\Delta t$ y después, si el modo adaptativo está activado, puede sustituir ese paso. Es un detalle de orden real que exige revisar balance con datos de ejecución; no se supone conservación exacta ni error cuantificado sin medirlo.

Cuando se habilita \texttt{adaptive\_ts}, la regla nativa de cada estado húmedo y borde es
\begin{equation}
 \Delta t\leq\min\!\left(\Delta t_{\mathrm{inicial}},
 C_{\mathrm{CFL}}\frac{\Delta x}{|u|+\sqrt{gh}},
 C_{\mathrm{CFL}}\frac{\Delta y}{|v|+\sqrt{gh}}\right).\label{eq:cfl}
\end{equation}
Las fracciones tienen segundos: m/(m/s). \texttt{dg2} fija por defecto $C_{\mathrm{CFL}}=0.33$, \texttt{DepthThresh}=0.001\,m, \texttt{DG2DepthThresh}=$-1$ y \texttt{DG2ThinDepthTstep}=0; un paso delgado positivo configurado puede imponer otra cota. \texttt{SpeedThresh} se inicializa en la fuente nativa para evitar operaciones con velocidades minúsculas. Los umbrales son mecanismos internos, \emph{no} tolerancias de error físico aprobadas. \src{lisflood.cpp} inicia \texttt{adaptive\_ts=ON}; el archivo \texttt{.par} de esta ruta declara \texttt{initial\_tstep}, pero no contiene opción que cambie explícitamente esa bandera. Por tanto, el camino normal usa recálculo CFL según ese default nativo, sujeto a comprobar la configuración efectiva de cada ejecución; no atribuir el modo adaptativo únicamente a \texttt{initial\_tstep}. El limitador de pendientes existe (Krivodonova/minmod en \src{swe/dg2/slope_limiter.cpp}) pero \texttt{limit\_slopes} inicia desactivado y la ruta ROI no lo activa; los bordes tienen tratamiento propio. Ningún umbral implica precisión, estabilidad o resolución de frente demostrada. Fuentes: \src{pars.cpp} (opción \texttt{dg2}), \src{lisflood.cpp} (inicialización), \src{swe/dg2.cpp} (\texttt{Tstep\_from\_cfl}, \texttt{min\_dt}, \texttt{solve}).

\section{Lluvia sintética: de la petición a un volumen}
\label{sec:lluvia}
La solicitud acepta cantidad acumulada $P$ y duración $T_r$ o intensidad y duración; \src{pipeline/hydrology/inputs.py} normaliza ambos modos a $P$ en m, $T_r$ en s y media $\bar r=P/T_r$ en m/s. Una intensidad expresada en mm/h se convierte con $r_{\mathrm{m/s}}=r_{\mathrm{mm/h}}/(1000\cdot3600)$. Cantidad en mm \emph{no} equivale a intensidad en mm/h hasta dividir por su intervalo temporal. La lluvia es definida por el usuario, no observada ni estimada a partir de radar.

Con número \emph{par} $N$ de intervalos iguales, $\Delta t_r=T_r/N$, el perfil \texttt{uniform\_balanced} genera $N/2$ desviaciones $d_j\sim\mathcal U[-\delta,\delta]$, sus opuestas $-d_j$, y reordena las $N$ contribuciones mediante semilla; aquí $\delta=\min(10/1000/3600\ \unit{m/s},\bar r)$. Luego $r_k=\bar r+d_k\geq0$ y
\begin{equation}
 \sum_{k=0}^{N-1}r_k\Delta t_r
 =N\bar r\frac{T_r}{N}+\Delta t_r\sum_kd_k=P.\label{eq:rainbal}
\end{equation}
Esta igualdad nominal conserva la lámina prescrita del \emph{perfil}; la integración real del solver con pasos propios y redondeos requiere comprobación separada. La cota solicitada $10$\,mm/h es límite de perturbación, no lluvia máxima: si la media es menor, la cota efectiva es menor; la lluvia instantánea puede superar la media. La misma tasa de cada intervalo se aplica a todos los \emph{centros de celda} dentro de la ROI; fuera es cero. La selección por centro es una discretización, no un campo meteorológico observado continuo. Si cada centro seleccionado tiene área $\Delta A$, el volumen nominal del intervalo es $V_k=r_k\Delta t_r N_{\mathrm{ROI}}\Delta A$, con m$^3$; para ROI recortada geométricamente no equiparar necesariamente $N_{\mathrm{ROI}}\Delta A$ a su área exacta. Véanse \src{pipeline/hydrology/rainfall.py} y \src{pipeline/hydrology/rainfall_netcdf.py}.

El exportador escribe NetCDF3 de desplazamiento de 64 bits, cubo \texttt{rainfall\_depth(time,y,x)} con \emph{mm acumulados por intervalo siguiente}, no m/s: $R_k=1000r_k\Delta t_r$. El eje \texttt{time} son horas desde cero, $x,y$ centros en metros, $x$ creciente y $y$ descendente; añade cuadro de lluvia nula al final y otro al término de la simulación si hay cola seca. El lector nativo \src{rain/rain.tpp} convierte el acumulado a tasa mediante
\begin{equation}
 r_k=\frac{R_k}{1000\,[3600(t_{k+1}^{\rm h}-t_k^{\rm h})]}\quad\unit{m/s},
 \qquad \Delta h=r_k\Delta t.\label{eq:ncrain}
\end{equation}
Para la contabilidad interna añade $\Delta V=\Delta h\Delta A$ por celda. El lector \src{lisflood2/file_tool.cpp} usa identificadores de variables como identificadores de dimensiones al consultar longitudes: el exportador establece deliberadamente el orden del encabezado y un eje ancla para compatibilidad con SciPy y \emph{esta versión} nativa. No generalizar el orden de variables a NetCDF arbitrarios ni atribuirle validez física. El lector usa horas por construcción, no interpreta libremente atributos alternativos de tiempo. La duración total debe cubrir la lluvia y su cola; la serie nula evita extender tácitamente la última tasa.

\section{Terreno, dominio, contorno e inicio}
\label{sec:entradas}
\subsection{De NASADEM a la malla de cálculo}
La ROI y \texttt{domain} llegan como polígonos EPSG:4326; el dominio debe cubrir la ROI. Tras proyectarlos a un CRS métrico válido elegido por la persona solicitante (por ejemplo, EPSG:32615 solo si es apropiado para el lugar), \src{pipeline/hydrology/roi_runner.py} alinea la caja de contorno a múltiplos de resolución, expande dos celdas por lado y obtiene un \emph{rectángulo completo}. No se simula solo el recorte poligonal: los bordes hidráulicos están en el rectángulo. El contorno inverso densificado determina teselas a adquirir. Se rechazan longitudes que crucen el antimeridiano y dominios fuera de latitudes NASADEM $-56^\circ$ a $60^\circ$; una tesela que simplemente intersecta la ROI no garantiza cobertura de toda la malla.

\src{pipeline/elevation/nasadem.py} adquiere HGT, \src{pipeline/hydrology/dem_prepared.py} prepara elevaciones y \texttt{roi\_runner} remuestrea la rejilla preparada a cada centro de la malla alineada. Se rechaza cualquier centro sin elevación finita; no se sustituye nodata por cero ni se inventa terreno. La operación vertical declarada conserva cotas EGM96: solo se reproyecta \emph{horizontalmente}. Los HGT de un segundo de arco son muestras angulares, no resolución métrica fija; el método \texttt{nearest} o \texttt{bilinear} solicitado no crea microtopografía. La idoneidad de cauces, diques naturales, vegetación, fecha y datum de niveles externos requiere estudio adicional. \texttt{prepare\_dg2\_dem} genera las tres capas modales \texttt{.dem}, \texttt{.dem1x}, \texttt{.dem1y}, con checks de geometría y hashes, sin certificar calidad topográfica.

\subsection{Frontera idealizada y condición inicial}
\texttt{dry\_zero\_velocity} implica $h(x,y,0)=q_x(x,y,0)=q_y(x,y,0)=0$ en el dominio: hipótesis impuesta, no prueba de que cauces, humedales o agua permanente estén secos. Sin \texttt{bcifile} el borde nativo por defecto es \texttt{NONE} (tratamiento reflectante de momento normal); con \texttt{boundary:free}, \src{pipeline/hydrology/boundaries.py} escribe \texttt{FREE} sobre los cuatro lados completos. En \src{swe/dg2.cpp}, \texttt{set\_boundary\_values}, \texttt{NONE} invierte la componente de momento normal fuera y \texttt{FREE} mantiene estado exterior igual a la media interior: \texttt{FREE} puede permitir entrada o salida según el estado; \emph{no} es la identificación observada de un desagüe ni una condición de nivel medida. La elección incorrecta del contorno puede dominar más que el orden numérico. La ruta ROI no acepta hidrograma de entrada, nivel observado ni batimetría como campos de su solicitud. Tampoco aplica infiltración: \texttt{infiltration:"unsupported"} reconoce una limitación, no un parámetro físico de pérdida igual a cero obtenido de medidas.

\section{Ciclo de vida operativo y API}
\label{sec:vida}
\begin{center}
\begin{tikzpicture}[node distance=3mm,box/.style={draw,rounded corners,fill=blue!6,align=center,text width=2.35cm,minimum height=1.1cm},arr/.style={-{Stealth},thick}]
\node[box] (r) {ROI/dominio\\JSON validado};
\node[box,right=of r] (d) {NASADEM\\DEM métrico};
\node[box,right=of d] (f) {Coeficientes DG2\\lluvia y frontera};
\node[box,right=of f] (s) {Motor DG2 CPU\\ASCII WD/MAX};
\node[box,right=of s] (c) {Control de salidas\\COG y API privada};
\draw[arr] (r)--(d);\draw[arr] (d)--(f);\draw[arr] (f)--(s);\draw[arr] (s)--(c);
\end{tikzpicture}
\end{center}

\subsection{Diagrama A: de la solicitud al mapa descargable}
\label{sec:flujo-api}
El recorrido siguiente separa la \emph{admisión} (1--3), la \emph{preparación} (4--6), el \emph{cálculo} (7) y la \emph{publicación} (8--9). Cada caja indica entrada $\to$ transformación $\to$ salida; las ramas laterales son fallos, no resultados parciales publicables. El diagrama B abre la caja 7.
\begin{center}
\begin{tikzpicture}[
 node distance=3mm,
 paso/.style={draw,rounded corners,fill=blue!6,align=left,text width=9.1cm,
              minimum height=1.08cm,inner sep=4pt,font=\small},
 fallo/.style={draw,rounded corners,fill=red!9,align=center,text width=2.8cm,
               inner sep=4pt,font=\small},
 flecha/.style={-{Stealth},thick}]
\node[paso] (a1) {\textbf{1. Entrada HTTP autenticada} $\to$ JSON de ROI, dominio, lluvia, terreno, hidráulica y tiempos $\to$ petición para \texttt{start/}.};
\node[paso,below=of a1] (a2) {\textbf{2. Validación previa} $\to$ geometría, dominio cubre ROI, CRS métrico, hipótesis y presupuestos $\to$ solicitud normalizada.};
\node[paso,below=of a2] (a3) {\textbf{3. Aceptación} $\to$ UUID y directorio privado $\to$ respuesta \texttt{202} y worker fuera del request.};
\node[paso,below=of a3] (a4) {\textbf{4. Adquisición} $\to$ teselas NASADEM HGT (cotas EGM96) $\to$ datos del dominio; no son profundidades.};
\node[paso,below=of a4] (a5) {\textbf{5. Terreno} $\to$ rectángulo métrico completo y DEM finito $\to$ \texttt{.dem/.dem1x/.dem1y}, ec.~\eqref{eq:vertices}.};
\node[paso,below=of a5] (a6) {\textbf{6. Forzamientos} $\to$ lluvia pareada y NetCDF, ecs.~\eqref{eq:rainbal}--\eqref{eq:ncrain}; borde \texttt{closed/free}.};
\node[paso,below=of a6] (a7) {\textbf{7. Cálculo} $\to$ parámetros e insumos con hashes $\to$ proceso DG2 CPU (véase diagrama B).};
\node[paso,below=of a7] (a8) {\textbf{8. Comprobación} $\to$ ASCII \texttt{.wd/.max} $\to$ malla, nodata, profundidad no negativa y cuadros ordenados.};
\node[paso,below=of a8] (a9) {\textbf{9. Publicación privada} $\to$ COG en CRS métrico $\to$ \texttt{status/results/artifacts} para propietario.};
\foreach \prim/\sig in {a1/a2,a2/a3,a3/a4,a4/a5,a5/a6,a6/a7,a7/a8,a8/a9}
  {\draw[flecha] (\prim.south)--(\sig.north);}
\node[fallo,right=7mm of a2] (rechazo) {\textbf{No admitido}\\400 solicitud o 503 servicio\\sin trabajo creado};
\draw[flecha] (a2.east)--node[above,font=\scriptsize]{error} (rechazo.west);
\node[fallo,right=7mm of a7] (fallo) {\textbf{Trabajo fallido}\\\texttt{failed} o\\\texttt{timed\_out}\\sin resultado disponible};
\draw[flecha] (a7.east)--node[above,font=\scriptsize]{error/plazo} (fallo.west);
\end{tikzpicture}
\end{center}
\noindent\textbf{Correspondencia de código.} 1--3: \src{backend/analyses/roi_flood_runs.py} (\texttt{RoiRunStartView}) y \src{pipeline/hydrology/roi_runner.py} (\src{validate_roi_run}, \src{start_roi_run}). 4--6: \src{run_roi_worker/_prepare} del mismo módulo, \src{pipeline/elevation/nasadem.py}, \src{pipeline/hydrology/dem_dg2.py}, \src{pipeline/hydrology/rainfall.py}, \src{pipeline/hydrology/rainfall_netcdf.py} y \src{pipeline/hydrology/boundaries.py}. 7: \texttt{run\_roi\_worker} crea \texttt{roi.par} y lanza el ejecutable fijado en servidor. 8--9: \src{pipeline/hydrology/results.py} (\texttt{summarize\_results}, \texttt{export\_depth\_cog}) y las vistas privadas del backend. Una falla de adquisición o preparación, un plazo agotado o una salida inválida terminan como \texttt{failed/timed\_out}; la flecha lateral de 7 resume fallos de \emph{cualquier etapa del worker}, no solo del cálculo. La admisión 202 no acredita éxito ni validación física.

\newpage
\subsection{Diagrama B: un paso del bucle nativo DG2}
\label{sec:flujo-nativo}
La ecuación \eqref{eq:conserv} gobierna el \emph{modelo continuo}; el diagrama muestra el \emph{orden del programa} en \src{data/tools/lisflood-fp-8.0/runtime/source/swe/dg2.cpp}, función \texttt{dg2::solve}. Aquí $\Delta t$ puede cambiar dentro de la vuelta: no igualar la lluvia aplicada al comienzo con la duración CFL recalculada después sin comprobar el balance.
\begin{center}
\begin{tikzpicture}[
 node distance=2.6mm,
 paso/.style={draw,rounded corners,fill=blue!6,align=left,text width=9.05cm,
              minimum height=.99cm,inner sep=3pt,font=\small},
 decision/.style={paso,fill=orange!12},
 final/.style={draw,rounded corners,fill=green!9,align=center,text width=2.95cm,
               inner sep=3pt,font=\small},
 flecha/.style={-{Stealth},thick}]
\node[paso] (b1) {\textbf{1. Inicio} $\to$ DEM modal, $h=q_x=q_y=0$ (supuesto seco), lluvia NetCDF y condiciones de borde.};
\node[decision,below=of b1] (b2) {\textbf{2. ¿Continúa?} $t<\texttt{Sim\_Time}$: si sí, ejecutar la vuelta; si no, escribir máximo y salir.};
\node[paso,below=of b2] (b3) {\textbf{3. Lluvia} $\to$ \texttt{rain.update\_H}: mm/intervalo $\to$ m/s $\to$ $h{+}=r\Delta t$ \emph{vigente}, ec.~\eqref{eq:ncrain}.};
\node[paso,below=of b3] (b4) {\textbf{4. Intercambios y bordes} $\to$ actualizar fuentes puntuales nativas y valores de frontera; la API ROI no configura aportes puntuales.};
\node[paso,below=of b4] (b5) {\textbf{5. Paso temporal} $\to$ con modo adaptativo, recalcular $\Delta t$ mediante CFL, ec.~\eqref{eq:cfl}.};
\node[paso,below=of b5] (b6) {\textbf{6. Fricción} $\to$ \texttt{apply\_friction} fuera de RK, actualización de momentos, ec.~\eqref{eq:fricdisc}.};
\node[paso,below=of b6] (b7) {\textbf{7. Dos etapas RK} $\to$ reconstrucción de caras, $h^*$, HLL y residuos $L_0,L_{1x},L_{1y}$ en \emph{cada} etapa; ecs.~\eqref{eq:weak}, \eqref{eq:hll}, \eqref{eq:rk}.};
\node[paso,below=of b7] (b8) {\textbf{8. Diagnóstico} $\to$ actualizar máximo, acumular estadísticas de flujo y luego $t{+}=\Delta t$ e iteración.};
\node[paso,below=of b8] (b9) {\textbf{9. Guardado condicional} $\to$ estadísticas de masa y cuadros cuando $t$ rebasa umbrales $\to$ volver a 2.};
\foreach \prim/\sig in {b1/b2,b2/b3,b3/b4,b4/b5,b5/b6,b6/b7,b7/b8,b8/b9}
  {\draw[flecha] (\prim.south)--(\sig.north);}
\node[final,right=7mm of b2] (fin) {\textbf{No}\\\texttt{write\_max\_field}\\retorno al diagrama A};
\draw[flecha] (b2.east)--(fin.west);
\draw[flecha] (b9.west)--++(-7mm,0) |- (b2.west);
\end{tikzpicture}
\end{center}
\noindent\textbf{Puntos de auditoría.} \texttt{rain.update\_H} en \src{rain/rain.tpp} antecede a \texttt{Tstep\_from\_cfl}; \texttt{apply\_friction} está fuera de \texttt{rk\_stage1/2}. Cada etapa usa \src{swe/dg2/modifiedvars.cpp} para caras hidroestáticas, \src{swe/hll_include.h} para el flujo y \texttt{L0/L1x/L1y} de \src{swe/dg2.cpp} para las tres amplitudes por campo (véanse también \eqref{eq:basis} y \eqref{eq:mean}). El limitador está disponible, \emph{no activado por defecto}; no hay llamada a infiltración en este bucle. El guardado ocurre tras avanzar $t$ sin forzar pasos para coincidir con \texttt{saveint} o el fin: el ASCII no guarda el tiempo físico de cada cuadro, y la API comunica \texttt{time\_s:null}. Estas etapas no incluyen una validación de precisión hidráulica.
\newpage
\noindent Un \texttt{POST /api/flood-scenarios/roi/start/} autenticado valida en memoria el JSON, polígonos, CRS, cobertura geográfica teórica y límites antes de crear trabajo o iniciar adquisición. Solo acepta siete secciones exactas: \texttt{roi}, \texttt{domain}, \texttt{rainfall}, \texttt{rainfall\_profile}, \texttt{terrain}, \texttt{hydraulics}, \texttt{simulation}. Los archivos y binarios son definidos por el servidor, no por el cliente. El POST devuelve 202 con UUID, \texttt{status\_url} y \texttt{results\_url}; el worker fuera del request gestiona adquisición y ejecución. No se afirma cola Celery ni cuota automática de CPU. El estado se conserva en un directorio privado por UUID/propietario con escritura atómica de \texttt{state.json}; cache NASADEM por trabajo y hashes de entradas, ejecutable y utilidad, parámetro, lluvia y COG. Los grupos de proceso de preparación y el solver tienen plazos separados. Un subproceso no constituye por sí solo un sandbox de seguridad.

\begin{center}
\begin{tikzpicture}[node distance=5mm and 5mm,box/.style={draw,rounded corners,fill=blue!6,align=center,text width=1.83cm,minimum height=.95cm,font=\small},bad/.style={box,fill=red!8},arr/.style={-{Stealth},thick}]
\node[box] (a) {created};\node[box,right=of a] (b) {acquiring};\node[box,right=of b] (c) {preparing};\node[box,right=of c] (d) {running};\node[box,right=of d] (e) {results};\node[box,right=of e] (f) {succeeded};
\foreach \u/\v in {a/b,b/c,c/d,d/e,e/f}{\draw[arr] (\u)--(\v);}
\node[bad,below=1cm of c] (g) {failed / timed\_out};
\draw[arr] (b)--(g);\draw[arr] (d)--(g);\draw[arr] (e)--(g);
\end{tikzpicture}
\end{center}
\noindent\texttt{GET /api/flood-roi-runs/<uuid>/} expone estado; fallos devuelven códigos seguros sin traceback ni rutas internas. \texttt{GET .../results/} responde 409 hasta éxito; en éxito incluye geometría de malla y CRS, unidades \texttt{m above local bed}, cuadros ordenados, estadísticas, máximo, advertencias y metadatos de artefactos. \texttt{GET .../artifacts/<nombre>/} solo sirve TIFF listado de trabajo exitoso al propietario y comprueba SHA256; respuesta completa con \texttt{ETag}, \emph{sin HTTP Range}. Acceso ajeno o nombre no autorizado produce 404. No se expone el ASCII original por esa API. Véanse \src{backend/analyses/roi_flood_runs.py} y \src{docs/api-simulacion-roi.md}.

\subsection{Solicitud completa de ejemplo, no caso calibrado}
\begin{small}\begin{verbatim}
{
 "roi":{"type":"Polygon","coordinates":[[[-92.9998,16.8],[-92.9995,16.8],
 [-92.9995,16.8003],[-92.9998,16.8003],[-92.9998,16.8]]]},
 "domain":{"type":"Polygon","coordinates":[[[-92.9999,16.7998],
 [-92.9993,16.7998],[-92.9993,16.8005],[-92.9999,16.8005],
 [-92.9999,16.7998]]]},
 "rainfall":{"amount":2,"unit":"mm","duration":1,"duration_unit":"s"},
 "rainfall_profile":{"interval_count":2,"seed":7,
                     "distribution":"uniform_balanced"},
 "terrain":{"target_crs":"EPSG:32615","resolution_m":20,
            "resampling":"nearest"},
 "hydraulics":{"manning":0.02,"initial_condition":"dry_zero_velocity",
               "boundary":"free","infiltration":"unsupported",
               "domain_nonurban":true},
 "simulation":{"duration_s":2,"save_interval_s":1,"initial_tstep_s":0.1}
}
\end{verbatim}\end{small}
Todos los números de este ejemplo ilustran un contrato, no constituyen tormenta probable, rugosidad recomendada, precisión de 20\,m ni recomendación de duración. La lluvia del ejemplo cae durante 1\,s y la simulación dura 2\,s: la cola restante es seca. La tasa media nominal sería $0.002$\,m/s ($7200$\,mm/h), matemáticamente coherente con 2\,mm en 1\,s pero físicamente extrema; por eso \emph{no} puede presentarse como evento representativo. La selección de dominio y supuestos exige intervención experta.

\noindent\begin{tabularx}{\linewidth}{p{.23\linewidth}X}
\toprule Respuesta & Interpretación correcta\\\midrule
\texttt{202 id/state/urls} & Identifica trabajo admitido; no indica simulación terminada ni validada.\\
\src{frames[i].frame_index} & Orden para reproducir cuadros. \texttt{time\_s:null} significa tiempo físico de ese ASCII no registrado. \texttt{scheduled\_time\_s=i*save\_interval\_s} es tiempo solicitado, no medido.\\
\texttt{maximum} & Máximo de profundidad provisto por el motor; no corresponde a una hora de llegada ni necesariamente a un único cuadro.\\
\texttt{grid/nodata} & CRS métrico y centinela del raster; nodata es desconocido, no profundidad cero.\\
\texttt{artifacts} & TIFF COG numérico con SHA256 y límites geográficos para encuadre, no imagen coloreada Web Mercator ni URL de teselas XYZ.\\
\src{synthetic:true}, \src{hydraulic_accuracy_validated:false} & Identifican carácter hipotético y ausencia de certificación física; no se suprimen en éxito.\\\bottomrule
\end{tabularx}

El motor comprueba guardado \emph{después} de actualizar $t$ y no trunca los pasos al próximo \texttt{saveint} ni al \texttt{sim\_time}. Por eso el adaptador \src{pipeline/hydrology/results.py} ordena índices y calcula solo horas \emph{programadas}; nunca presenta \texttt{scheduled\_time\_s} como observación física exacta. Los ASCII \texttt{.wd} y \texttt{.max} se validan por encabezado, forma, finitud, no negatividad fuera de nodata y malla común; después se exportan a COG \texttt{float32} en CRS proyectado, con hashes origen/destino. La estadística \texttt{wet\_cells} cuenta valores válidos $>0$, \emph{no} una extensión oficial con umbral hidrológico validado. Para MapLibre hay que decodificar/reproyectar a Web Mercator o montar teselas autenticadas por separado: entregar coordenadas EPSG:4326 de encuadre no convierte los píxeles UTM ni agrega soporte Range. No existe API de velocidad ni tiempo de llegada.

\section{Parámetros de servicio, trazabilidad y errores}
\label{sec:operacion}
El archivo \texttt{roi.par} generado declara \texttt{dg2}, \src{DEMfile} \src{terrain.dem}, \texttt{sim\_time}, \texttt{saveint}, \src{initial_tstep}, \src{fpfric} y \src{dynamicrainfile} \src{rain.nc}; para \texttt{free} añade \texttt{bcifile}. La malla de datos y lluvia, la utilidad DEM y el ejecutable se comprueban frente a rutas y hashes de configuración. SHA256 de binario \src{777d50f8d6a926260baa31b4c8ca3ecc0b3ae710494bebb2d4699ccf8c39775c}; de \texttt{generateDG2DEM} \src{ada476fc141d536d09abd4aa1238d192bbddfd627574886ada12784ba0e964a8}. Los hashes identifican bytes usados, no prueban exactitud física ni una licencia de redistribución. Release fuente LISFLOOD-FP 8.0, ejecutable local identificado como 8.0.3 float CPU; instantánea de repositorio al solicitar este documento: \src{615d0ce23eee51df7c130c93a1306958b2d0f0c6}. El \texttt{CODE\_VERSION} configurable identifica la versión declarada del adaptador y debe mantenerse actualizado por operación: su valor por defecto no es garantía de correspondencia automática con HEAD.

\noindent\begin{tabularx}{\linewidth}{p{.28\linewidth}p{.15\linewidth}X}
\toprule Ajuste inicial del servidor & Valor & Significado y restricción\\\midrule
\texttt{FLOOD\_ROI\_DG2\_ENABLED} & true & \texttt{false} desactiva \texttt{start/}; no acredita acceso Earthdata.\\
\texttt{MAX\_CELLS / MAX\_RAIN\_VALUES} & 1000 / 5000 & Presupuestos de malla y valores, no resolución física requerida ni tolerancia hidráulica.\\
\texttt{MAX\_FRAMES / MAX\_TILES} & 5 / 2 & Presupuestos de salidas y teselas; el área contribuyente puede necesitar más.\\
\texttt{TIMEOUT\_S / PREPARATION\_TIMEOUT\_S} & 15 / 120\,s & Plazos del solver y de la fase preparatoria, no duración simulada.\\\bottomrule
\end{tabularx}
El comando \texttt{PYTHONPATH=. python backend/manage.py check\_flood\_service} verifica preparación local sin red; \texttt{--earthdata}, solo con autorización del operador, comprueba una adquisición distinta. No son tests de física. La lectura de credenciales ocurre en entorno privado, nunca en JSON ni en los artefactos públicos. Una solicitud mal formada o fuera de presupuesto es rechazada antes de crear trabajo (400); servicio no disponible o configuración inválida (503); no autenticado (401/403); recurso ajeno o desconocido (404); resultados aún no listos (409). La fase de adquisición puede fallar por credenciales/cobertura; preparación puede fallar por nodata, formato o plazo; el solver por timeout/salida no válida. Los códigos expuestos son sanitizados y no transmiten excepciones ni rutas privadas. \src{backend/config/flood_service.py}, \src{pipeline/hydrology/roi_runner.py}, \src{backend/analyses/roi_flood_runs.py} son la correspondencia operativa.

\section{Evidencia observada y lo que no demuestra}
\label{sec:evidencia}
Se conservan \emph{dos evidencias distintas}, que no deben fusionarse en un resultado regional real:
\begin{enumerate}
\item La unión de adaptador y motor real con proveedor \emph{HGT sintético} terminó en \texttt{succeeded}, malla $9\times8$ a 20\,m, tres cuadros WD y un máximo, cuatro COG EPSG:32615. Registro: \src{data/test-runs/roi-dg2-136934da-7370-4113-b337-36ebc0bcafcc/6c5b38f4-9d14-4c71-912b-04ed7629d94a/state.json}. Prueba continuidad técnica de esa ejecución concreta, no exactitud ni disponibilidad Earthdata para el mismo run.
\item La adquisición \emph{real Earthdata}, separada, guardó comprobante de tesela NASADEM \texttt{n16w094}, ZIP de 14\,251\,621 bytes: \src{data/service-checks/7507c3c0-8ba1-43b1-920e-a35fa58ea70f/proof.json}. Prueba descarga de esa tesela en esa comprobación, no un escenario DG2 ejecutado sobre ella ni autenticidad criptográfica independiente del servidor remoto.
\end{enumerate}
\begin{tabularx}{\linewidth}{p{.23\linewidth}p{.37\linewidth}X}
\toprule Nivel & Evidencia disponible & Comprobación todavía necesaria\\\midrule
Formato e integración & Manifiestos y hashes, NetCDF nativo, ASCII/COG de seam sintético. & Repetir por caso y comparar entradas/salidas; no extrapolar a dominio grande.\\
Matemática numérica & Código de operadores y dimensionalidad descritos aquí. & Ensayos lago en reposo, lecho variable, mojado/seco, convergencia de malla y tiempo.\\
Balance volumétrico & Registro de lluvia y archivo nativo \texttt{.mass}. & Investigar \texttt{Qerror/Verror}, extraer flujos, contrastar $V_f-V_i$ con aportes y fronteras con tolerancia justificada antes de afirmar conservación.\\
Fidelidad física & Ningún evento regional validado en estas evidencias. & DEM/cauces/batimetría, lluvia y niveles independientes, Manning, infiltración, sensibilidad y contraste espacio-temporal externo.\\\bottomrule
\end{tabularx}
Como control, para cada paso ideal $V_{n+1}-V_n\stackrel{?}{=}\Delta t\sum_Kr_K\Delta A-\Delta t\sum_{f\subset\partial\Omega}\widehat q_f\ell_f$ (m$^3$), con convención de normal saliente y tiempo realmente empleado. No usar esta igualdad como resultado observado: hay que determinar cómo el motor integra y registra los aportes y sobrepasos temporales. \texttt{getfattr} registró una advertencia del smoke pendiente de investigación; no se ocultó. Los resultados siguen marcados \texttt{hydraulic\_accuracy\_validated:false}. Un mapa raster que visualmente parece inundación no mide precisión, recall, área verdadera ni error de profundidad. Las escenas NISAR pueden servir de contraste de extensión con hora y máscaras verificadas; una diferencia de retrodispersión no demuestra agua, profundidad o causalidad por sí sola. Sin referencia independiente, el escenario es exploratorio.

\section{Mapa de ecuaciones a código y puntos de auditoría}
\label{sec:mapa}
En la tabla, las rutas nativas se entienden bajo \src{data/tools/lisflood-fp-8.0/runtime/source/}; los nombres de función son anclas auditables más robustas que números de línea cuando cambie el release.
\begin{longtable}{p{.19\linewidth}p{.34\linewidth}p{.38\linewidth}}
\toprule Relación y unidad & Función y salida & Control de interpretación\\\midrule\endhead
\eqref{eq:conserv}--\eqref{eq:flux}, m/s y m$^2$/s$^2$ & \src{swe/flux.cpp}: \texttt{physical\_flux\_x/y}; \src{swe/dg2.cpp}: \texttt{L0/L1x/L1y} & No confundir residuo modal con solo una media; derivadas dividen por metros.\\
\eqref{eq:vertices}, m & \src{preprocess/generateDG2DEM.cpp} y \src{swe/dg2/fields.cpp}: tres \texttt{.dem*} & Comparar cotas y centinela; los coeficientes no son pendientes GIS en m/m.\\
$z^*,h^*,q^*$, m, m, m$^2$/s & \src{swe/dg2/modifiedvars.cpp}: \texttt{initialise\_Zstar}, \texttt{update\_Hstar}, \texttt{discharge\_star} & Cara seca no implica prueba de positividad global.\\
\eqref{eq:hll}, m$^2$/s y m$^3$/s$^2$ & \src{swe/hll_include.h}: rama intermedia; \src{swe/dg2.cpp}: \texttt{update\_fluxes} & Componente transversal usa contacto elegido, no HLL vectorial directo.\\
\eqref{eq:fricdisc}, m$^2$/s & \src{swe/dg2/friction.cpp}: \texttt{friction}, \texttt{apply\_friction} & Separar pendiente adimensional y fuerza de momento; operador fuera RK.\\
\eqref{eq:cfl}, s; \eqref{eq:rk}, unidades de $U$ & \src{swe/dg2.cpp}: \texttt{min\_dt}, \texttt{rk\_stage1/2}, \texttt{solve} & Paso CFL condicional a modo adaptativo; no sincroniza exactamente cuadros.\\
\eqref{eq:rainbal}--\eqref{eq:ncrain}, m y m/s & \src{pipeline/hydrology/rainfall.py}, \src{pipeline/hydrology/rainfall_netcdf.py}, \src{rain/rain.tpp} & Diferenciar acumulado mm, horas de archivo y m/s del motor; revisar orden previo al CFL.\\
Profundidad $h$, m sobre lecho & \src{pipeline/hydrology/results.py}: \texttt{summarize\_results}, \texttt{export\_depth\_cog} & Nodata desconocido; tiempo real ausente; COG numérico en CRS métrico.\\\bottomrule
\end{longtable}

\section{Límites de interpretación y tareas para un caso científico}
\label{sec:limites}
Antes de interpretar una simulación real deben fijarse evento, dominio contribuyente, origen temporal UTC, calidad y resolución del DEM, datum de niveles, batimetría/cauces, flujo inicial, condición observada de borde, lluvia representativa, infiltración (no soportada en esta ruta), Manning espacial o justificación de valor uniforme y relación con zonas urbanas omitidas. Una frontera FREE no resuelve automáticamente estas preguntas. Establecer umbral de extensión y máscaras de nodata y agua permanente; comparar una observación al tiempo que efectivamente representa el cuadro, no el máximo temporal frente a una instantánea. Mientras el ASCII no registre tiempo físico, esa comparación puntual debe declarar su incertidumbre temporal o no realizarse. Elegir y registrar tolerancias de balance y convergencia \emph{antes} de usar resultados; no inventar porcentajes de error aceptables. Comparar con datos independientes separados de cualquier ajuste. El coste de una malla mayor se debe medir: los límites iniciales de servidor son presupuestos operativos pequeños, no capacidad regional verificada.

\section{Fuentes y correspondencia de versión}
\begin{enumerate}
\item Shaw, J. y cols. (2021), ``LISFLOOD-FP 8.0: the new discontinuous Galerkin shallow-water solver for multi-core CPUs and GPUs'', \emph{Geoscientific Model Development} 14, 3577--3602. \url{https://doi.org/10.5194/gmd-14-3577-2021}. Fuente de contexto del método publicado, no sustituto del cotejo del ejecutable concreto.
\item LISFLOOD-FP 8.0, artefacto original: \url{https://zenodo.org/records/4073011}; guía de los desarrolladores: \url{https://www.seamlesswave.com/LISFLOOD8.0}. Las ecuaciones atribuidas a \emph{esta ejecución} se cotejan con la fuente nativa local bajo \src{data/tools/lisflood-fp-8.0/runtime/source/}, no se atribuyen a una extracción incompleta de la página HTML del artículo.
\item NASADEM User Guide V1.3, producto HGT y referencia vertical EGM96: \url{https://lpdaac.usgs.gov/documents/2237/NASADEM_User_Guide_V13.pdf}. La referencia no certifica una elevación particular.
\item Contratos actuales: \src{docs/api-simulacion-roi.md}, \src{docs/configuracion-servicio-inundacion.md}, \src{docs/distribucion-lluvia.md}, \src{docs/dem-dg2.md}, \src{docs/resultados-dg2.md} y código citado en el \S\ref{sec:mapa}. Documentos conceptuales anteriores describen propuestas históricas y no desplazan el contrato implementado. La identificación de licencia en Zenodo no autoriza por sí sola redistribución del binario; revisar obligaciones antes de publicarlo.
\end{enumerate}
\end{document}
